% Options for packages loaded elsewhere
% Options for packages loaded elsewhere
\PassOptionsToPackage{unicode}{hyperref}
\PassOptionsToPackage{hyphens}{url}
\PassOptionsToPackage{dvipsnames,svgnames,x11names}{xcolor}
%
\documentclass[
  12pt,
]{article}
\usepackage{xcolor}
\usepackage[left=40mm,right=20mm]{geometry}
\usepackage{amsmath,amssymb}
\setcounter{secnumdepth}{5}
\usepackage{iftex}
\ifPDFTeX
  \usepackage[T1]{fontenc}
  \usepackage[utf8]{inputenc}
  \usepackage{textcomp} % provide euro and other symbols
\else % if luatex or xetex
  \usepackage{unicode-math} % this also loads fontspec
  \defaultfontfeatures{Scale=MatchLowercase}
  \defaultfontfeatures[\rmfamily]{Ligatures=TeX,Scale=1}
\fi
\usepackage{lmodern}
\ifPDFTeX\else
  % xetex/luatex font selection
  \setmainfont[]{Linux Libertine O}
\fi
% Use upquote if available, for straight quotes in verbatim environments
\IfFileExists{upquote.sty}{\usepackage{upquote}}{}
\IfFileExists{microtype.sty}{% use microtype if available
  \usepackage[]{microtype}
  \UseMicrotypeSet[protrusion]{basicmath} % disable protrusion for tt fonts
}{}
\usepackage{setspace}
% Make \paragraph and \subparagraph free-standing
\makeatletter
\ifx\paragraph\undefined\else
  \let\oldparagraph\paragraph
  \renewcommand{\paragraph}{
    \@ifstar
      \xxxParagraphStar
      \xxxParagraphNoStar
  }
  \newcommand{\xxxParagraphStar}[1]{\oldparagraph*{#1}\mbox{}}
  \newcommand{\xxxParagraphNoStar}[1]{\oldparagraph{#1}\mbox{}}
\fi
\ifx\subparagraph\undefined\else
  \let\oldsubparagraph\subparagraph
  \renewcommand{\subparagraph}{
    \@ifstar
      \xxxSubParagraphStar
      \xxxSubParagraphNoStar
  }
  \newcommand{\xxxSubParagraphStar}[1]{\oldsubparagraph*{#1}\mbox{}}
  \newcommand{\xxxSubParagraphNoStar}[1]{\oldsubparagraph{#1}\mbox{}}
\fi
\makeatother


\usepackage{longtable,booktabs,array}
\usepackage{calc} % for calculating minipage widths
% Correct order of tables after \paragraph or \subparagraph
\usepackage{etoolbox}
\makeatletter
\patchcmd\longtable{\par}{\if@noskipsec\mbox{}\fi\par}{}{}
\makeatother
% Allow footnotes in longtable head/foot
\IfFileExists{footnotehyper.sty}{\usepackage{footnotehyper}}{\usepackage{footnote}}
\makesavenoteenv{longtable}
\usepackage{graphicx}
\makeatletter
\newsavebox\pandoc@box
\newcommand*\pandocbounded[1]{% scales image to fit in text height/width
  \sbox\pandoc@box{#1}%
  \Gscale@div\@tempa{\textheight}{\dimexpr\ht\pandoc@box+\dp\pandoc@box\relax}%
  \Gscale@div\@tempb{\linewidth}{\wd\pandoc@box}%
  \ifdim\@tempb\p@<\@tempa\p@\let\@tempa\@tempb\fi% select the smaller of both
  \ifdim\@tempa\p@<\p@\scalebox{\@tempa}{\usebox\pandoc@box}%
  \else\usebox{\pandoc@box}%
  \fi%
}
% Set default figure placement to htbp
\def\fps@figure{htbp}
\makeatother


% definitions for citeproc citations
\NewDocumentCommand\citeproctext{}{}
\NewDocumentCommand\citeproc{mm}{%
  \begingroup\def\citeproctext{#2}\cite{#1}\endgroup}
\makeatletter
 % allow citations to break across lines
 \let\@cite@ofmt\@firstofone
 % avoid brackets around text for \cite:
 \def\@biblabel#1{}
 \def\@cite#1#2{{#1\if@tempswa , #2\fi}}
\makeatother
\newlength{\cslhangindent}
\setlength{\cslhangindent}{1.5em}
\newlength{\csllabelwidth}
\setlength{\csllabelwidth}{3em}
\newenvironment{CSLReferences}[2] % #1 hanging-indent, #2 entry-spacing
 {\begin{list}{}{%
  \setlength{\itemindent}{0pt}
  \setlength{\leftmargin}{0pt}
  \setlength{\parsep}{0pt}
  % turn on hanging indent if param 1 is 1
  \ifodd #1
   \setlength{\leftmargin}{\cslhangindent}
   \setlength{\itemindent}{-1\cslhangindent}
  \fi
  % set entry spacing
  \setlength{\itemsep}{#2\baselineskip}}}
 {\end{list}}
\usepackage{calc}
\newcommand{\CSLBlock}[1]{\hfill\break\parbox[t]{\linewidth}{\strut\ignorespaces#1\strut}}
\newcommand{\CSLLeftMargin}[1]{\parbox[t]{\csllabelwidth}{\strut#1\strut}}
\newcommand{\CSLRightInline}[1]{\parbox[t]{\linewidth - \csllabelwidth}{\strut#1\strut}}
\newcommand{\CSLIndent}[1]{\hspace{\cslhangindent}#1}



\setlength{\emergencystretch}{3em} % prevent overfull lines

\providecommand{\tightlist}{%
  \setlength{\itemsep}{0pt}\setlength{\parskip}{0pt}}



 


\usepackage{lineno}
\usepackage[noblocks]{authblk}
\usepackage{indentfirst}
\setlength{\parindent}{2em}
\renewcommand*{\Authsep}{, }
\renewcommand*{\Authand}{, }
\renewcommand*{\Authands}{, }
\renewcommand\Affilfont{\small}
\makeatletter
\newcommand*\correspondingauthor[1]{%
\g@addto@macro\@title{\par\normalsize\textsuperscript{*}Corresponding author: \texttt{#1}\par}
}
\correspondingauthor{m.klemp@dshs-koeln.de}
\usepackage{lua-ul}
\AtBeginDocument{%
  \renewcommand{\criticmarkupadd}[1]{{\color[HTML]{1A7321}#1}}%
  \renewcommand{\criticmarkuprm}[1]{{\color[HTML]{C62828}\strikeThrough{#1}}}%
  \renewcommand{\criticmarkupmark}[1]{{\color[HTML]{E65100}#1}}%
}
\makeatletter
\@ifpackageloaded{caption}{}{\usepackage{caption}}
\AtBeginDocument{%
\ifdefined\contentsname
  \renewcommand*\contentsname{Table of contents}
\else
  \newcommand\contentsname{Table of contents}
\fi
\ifdefined\listfigurename
  \renewcommand*\listfigurename{List of Figures}
\else
  \newcommand\listfigurename{List of Figures}
\fi
\ifdefined\listtablename
  \renewcommand*\listtablename{List of Tables}
\else
  \newcommand\listtablename{List of Tables}
\fi
\ifdefined\figurename
  \renewcommand*\figurename{Figure}
\else
  \newcommand\figurename{Figure}
\fi
\ifdefined\tablename
  \renewcommand*\tablename{Table}
\else
  \newcommand\tablename{Table}
\fi
}
\@ifpackageloaded{float}{}{\usepackage{float}}
\floatstyle{ruled}
\@ifundefined{c@chapter}{\newfloat{codelisting}{h}{lop}}{\newfloat{codelisting}{h}{lop}[chapter]}
\floatname{codelisting}{Listing}
\newcommand*\listoflistings{\listof{codelisting}{List of Listings}}
\makeatother
\makeatletter
\makeatother
\makeatletter
\@ifpackageloaded{caption}{}{\usepackage{caption}}
\@ifpackageloaded{subcaption}{}{\usepackage{subcaption}}
\makeatother
\makeatletter
\newcommand{\criticmarkupadd}[1]{}
\newcommand{\criticmarkuprm}[1]{#1}
\newcommand{\criticmarkupmark}[1]{#1}
\newcommand{\criticmarkupcomm}[1]{}
\makeatother
\usepackage{bookmark}
\IfFileExists{xurl.sty}{\usepackage{xurl}}{} % add URL line breaks if available
\urlstyle{same}
\hypersetup{
  pdftitle={Consensual Assessment Technique Methodology for Rating Creative Products: An Application to Football Actions},
  pdfauthor={Maximilian Klemp; Daniel Memmert; Robert Rein},
  pdfkeywords={consensual assessment technique, creativity
assessment, judge response theory, many-facet Rasch model, inter-rater
reliability, sport creativity},
  colorlinks=true,
  linkcolor={blue},
  filecolor={Maroon},
  citecolor={Blue},
  urlcolor={Blue},
  pdfcreator={LaTeX via pandoc}}


\title{Consensual Assessment Technique Methodology for Rating Creative
Products: An Application to Football Actions}


\author[1*]{Maximilian Klemp}
\author[1]{Daniel Memmert}
\author[1]{Robert Rein}

\affil[1]{Institute of Exercise Training and Sport Informatics, German
Sport University Cologne}
\affil[*]{Corresponding Author, m.klemp@dshs-koeln.de}


\date{}
\begin{document}
\maketitle
\begin{abstract}
The consensual assessment technique (CAT) is the dominant method for
assessing the creativity of products, yet its validity is
domain-contingent: the number and expertise of judges required cannot be
specified in general and must be established within each new domain of
application (Amabile, 1982; Cseh \& Jeffries, 2019). Sport is one such
domain in which this evidence has never been supplied. Using ratings of
discrete football actions as a test domain, this paper addresses how
many judges are needed for convergent creativity ratings, whether rater
expertise affects agreement or severity, whether rater-effects models
(Judge Response Theory and the Many-Facet Rasch Model) improve on simple
averaging, whether creativity is discriminable from aesthetic appeal and
technical correctness, and how long raters can perform the task
reliably. We propose a large-scale, multi-rater benchmark-and- pairwise
design in which approximately 1,000 to 2,000 purposively selected
actions will be rated on a 7-point unipolar novelty x appropriateness
scale by raters spanning an expertise gradient. Results are reported for
inter-rater reliability across judge counts, expertise comparisons,
JRT/MFRM rater diagnostics against naive averaging, and the
creativity/aesthetic/technical-correctness factor structure (all pending
data collection). The paper is intended as a reproducible methodological
template for CAT-based creative-product research, with football as the
applied testbed.
\end{abstract}

\newpage{}


\setstretch{1.2}
\section{Introduction}\label{introduction}

Creativity is widely regarded as a fundamental human competency, yet how
to measure it remains a live and unresolved methodological problem
across psychology, education, and applied fields -- not a solved
technical detail (Long et al., 2022). A taxonomy distinguishing the
creative person, process, product, and press (environment) has been
proposed precisely because each facet requires its own measurement logic
and cannot be assessed by a single all-purpose instrument (Barbot et
al., 2019; Batey, 2012). As the concluding piece of a recent special
issue on creativity assessment argues, clarifying which facet of
creativity is being measured is a necessary precondition for judging and
improving any specific method (Barbot et al., 2019; Barbot \&
Reiter-Palmon, 2019). The present paper is concerned with the
\emph{product} facet -- the creativity of a discrete, observable outcome
-- and with the subjective-judgment method that dominates its
assessment.

The consensual assessment technique (CAT) is the most widely used
approach to rating the creativity of products. Under CAT, judges with
relevant domain familiarity independently rate products' creativity
without being given a predefined rubric, on the logic that creativity is
best captured as a consensual, socially situated judgment rather than an
objectively specifiable property of the product itself (Amabile, 1982).
CAT has accumulated a substantial track record of reasonable inter-rater
agreement, particularly once ratings are averaged across multiple
judges, and treating rater as a facet of measurement error within a
generalizability-theory framework has clarified how much of that error
is attributable to raters specifically (Long \& Pang, 2015). Textbook
treatments now present CAT as a mature, practically usable tool
alongside divergent-thinking tests and self-report measures (Kaufman,
Plucker, et al., 2008).

Yet CAT's apparent maturity conceals a persistent, under-operationalized
gap. Cseh and Jeffries (2019) offer a systematic critique of the
technique, arguing that CAT's validity depends on the correct use of a
sufficient number of judges with an appropriate level of expertise, and
that both the required number and the required expertise are properties
of the specific product and domain being rated -- they cannot be fixed
once and for all across applications. On this view, every new domain in
which CAT is applied owes the field its own empirical evidence on these
design parameters. This paper's central premise is that such evidence
has not yet been supplied for the domain of sport, and that this
omission is representative of how CAT is often applied to new domains
more generally: with a design borrowed from precedent rather than
derived from domain-specific evidence.

\subsection{Sub-debate A: How many judges are
enough?}\label{sub-debate-a-how-many-judges-are-enough}

A persistent rule of thumb in the CAT literature holds that ``five to
ten judges'' is generally adequate for most purposes (Kaufman, Plucker,
et al., 2008). This heuristic is useful as a starting point but is not
itself derived from domain-specific empirical analysis. By contrast,
generalizability-theory approaches derive the number of judges (and
items) required for a target level of reliability directly from the data
at hand, and have shown that this number varies considerably by task
type and rating design -- in one divergent-thinking-rating context as
few as three untrained raters combined with four items sufficed for
adequate dependability, but the estimate did not generalize uniformly
across designs (Hass et al., 2018). The tension between a generic
heuristic and a domain-derived estimate is precisely the kind of gap
Cseh and Jeffries (2019) identify, and it motivates this paper's first
research question.

\subsection{Sub-debate B: Does rater expertise matter, and how
much?}\label{sub-debate-b-does-rater-expertise-matter-and-how-much}

A second, closely related debate concerns whether judges must be domain
experts. Amabile (1982) argues that judges' level of experience need not
be uniform, and that CAT's reliability holds up even when the products
being rated are naturalistically produced and non-parallel -- that is,
not generated under identical, standardized experimental conditions.
Direct empirical support for this claim comes from studies showing that
CAT inter-rater reliability remained high when expert judges rated
creative writing produced under a variety of different assignments
rather than a single controlled prompt (Baer et al., 2004), and from
design-engineering research modeling how novices and experts differ in
the knowledge they bring to evaluating creative work (Christiaans \&
Venselaar, 2005). At the same time, direct comparisons of expert and
non-expert raters using CAT show that non-experts can be meaningfully
discerning judges of creative products, complicating any simple
assumption that only experts should be recruited (Kaufman, Baer, et al.,
2008). Further evidence suggests that domain expertise does not
straightforwardly increase rater agreement: in a comparison of movie
ratings, raters who had simply seen more movies -- regardless of formal
expert status -- correlated more highly with professional critics than
less-exposed raters did (Plucker et al., 2009), suggesting that high
exposure or experience can substitute, at least partially, for formal
expertise. This ``heavy viewer'' finding directly motivates a design
choice pursued below: recruiting football raters who differ in exposure
and background rather than assuming only coaches or professionals can
serve as valid judges. Other work argues that domain-specific expertise
nonetheless matters for some purposes (Baer, 2015), and a further strand
of research cautions that expertise itself can introduce systematic
biases -- a ``dark side of expertise'' in which greater knowledge does
not uniformly translate into better judgment (Castel et al., 2007).
Rather than treating ``expert'' and ``novice'' as monolithic categories,
recent work has begun asking what individual-difference factors predict
discernment among novice raters specifically (Ceh et al., 2021). This
sub-debate is presented here as genuinely open, not settled, and is the
subject of this paper's second research question.

\subsection{Sub-debate C: How should ratings from multiple,
non-interchangeable judges be
aggregated?}\label{sub-debate-c-how-should-ratings-from-multiple-non-interchangeable-judges-be-aggregated}

A third, more technical debate concerns aggregation. The conventional
practice of averaging ratings across judges implicitly treats judges as
exchangeable random effects -- interchangeable sources of measurement
error that cancel out on average. This assumption ignores the
possibility that judges differ systematically in severity or leniency, a
concern that has motivated two related but historically separate
psychometric traditions. Judge Response Theory (JRT) proposes modeling
judgments with item-response- theory-style models in which judges, like
test items, have their own parameters (e.g., severity, discrimination),
yielding latent creativity estimates that are corrected for these
differences rather than averaged over them (Myszkowski \& Storme, 2019).
This proposal has since been operationalized in an open-source R package
that fits JRT models to CAT-style judgment data (Myszkowski, 2021) and
has been extended into a book-length treatment of item response theory
for creativity measurement more broadly (Myszkowski, 2024), with a
recent application decomposing novelty and appropriateness as separate
latent components within an item-response-tree framework (Myszkowski \&
Storme, 2025) -- a decomposition directly relevant to the
two-dimensional novelty x appropriateness scale adopted below.
Independently, the Many-Facet Rasch Model (MFRM) tradition, developed
originally for rater-mediated performance assessment, models rater
severity as one of several crossed facets alongside examinee ability and
item difficulty, and provides established diagnostics for detecting
systematic rater effects such as leniency, central tendency, and halo
(Eckes, 2015; Myford \& Wolfe, 2003, 2004). MFRM has been applied
directly to creativity ratings, showing that individual raters can be
diagnosed and examinees reliably differentiated despite idiosyncratic
rater behavior (Hung et al., 2012; Primi et al., 2019). Despite their
shared logic, JRT and MFRM emerged from separate literatures and have
not, to our knowledge, been directly compared against one another or
against naive averaging in a creativity-rating context at scale. This
paper's third research question addresses that gap.

\subsection{Construct validity: Is creativity discriminable from
adjacent
constructs?}\label{construct-validity-is-creativity-discriminable-from-adjacent-constructs}

A further concern, orthogonal to the number and expertise of judges, is
whether the construct being rated -- creativity -- is empirically
separable from adjacent constructs that judges might conflate with it,
such as aesthetic appeal or technical execution. Amabile (1982)
recommends that when CAT is applied to a new domain, judges should also
rate such secondary dimensions so that the creativity rating's
discriminant and construct validity can be assessed directly rather than
assumed. This recommendation is adopted here, using the standard
definitional consensus that creativity reflects the joint presence of
novelty (originality) and appropriateness (functionality, or fit to the
constraints and goals of the situation) (Runco \& Jaeger, 2012). This
paper's fourth research question asks whether creativity, aesthetic
appeal, and technical correctness form empirically distinct factors when
rated concurrently on the same set of products.

\subsection{A practical measurement concern: Rater
fatigue}\label{a-practical-measurement-concern-rater-fatigue}

Beyond these conceptual debates, a comparatively neglected practical
question is how long raters can perform a CAT-style rating task before
the quality or consistency of their judgments declines. This is a design
question relevant to any large-scale CAT study, not a sport-specific
one, and has been investigated directly in at least one prior study
examining whether CAT ratings remain reliable across extended rating
sessions (Barth \& Stadtmann, 2020). This paper's fifth research
question adapts that logic to the football domain.

\subsection{Narrowing to football: Why this domain, and why
now}\label{narrowing-to-football-why-this-domain-and-why-now}

Creativity is considered one of the key competencies to be developed and
evaluated in football players, and assessing it in an ecologically
valid, naturalistic environment is an explicit priority for the field
(Sa Fardilha \& Allen, 2019). Much of the existing sport-creativity
literature has studied the creative \emph{process} -- for example,
divergent-thinking ability and its development in players -- while the
creative \emph{product}, the discrete in-game action itself, has
received comparatively little attention (Memmert, 2010). One prior study
has evaluated the creativity of football actions directly, using a 0-10
expert rating scale, and demonstrated face validity by showing that
rated creativity predicted the temporal proximity of actions to goals
(Kempe \& Memmert, 2018); however, this study did not investigate how
many judges were needed, whether rater expertise mattered, or how
ratings from multiple judges should best be aggregated. Two further,
smaller-scale analogues exist: a two-expert-judge design rating in-game
creativity on a 1-10 scale achieved very high inter-judge correlation
(Oboeuf et al., 2020), and a training-transfer study used expert
on-field creativity ratings to compare functional-skills training
against divergent-thinking training (Zahno \& Hossner, 2022). Neither of
these studies systematically varied the number or expertise of judges,
nor did either apply JRT or MFRM to model rater severity. To our
knowledge, no CAT-in-sport study to date has systematically varied judge
number or expertise, or modeled rater severity using modern
rater-effects models, in this domain.

\subsection{Aim and research
questions}\label{aim-and-research-questions}

The aim of the present paper is to generate a large-scale,
systematically varied multi-rater creativity-rating dataset for a sports
domain (football actions) and to use it to provide the field with
empirically grounded, reproducible guidance on the psychometric design
choices that Cseh and Jeffries (2019) argue cannot be resolved in
general. Specifically, the paper addresses the following research
questions and, where the existing literature already supports a
directional prediction, states it as an explicit hypothesis rather than
leaving the question hypothesis-free by omission:

\begin{itemize}
\tightlist
\item
  \textbf{RQ1 (judges).} How many raters are needed for stable,
  convergent creativity ratings of football actions?

  \begin{itemize}
  \tightlist
  \item
    \textbf{H1.} Inter-rater reliability (ICC) will increase with the
    number of judges averaged into the aggregate rating, with
    asymptotically diminishing returns, consistent with
    generalizability-theory/ Spearman-Brown logic and with the empirical
    template set by Hass et al. (2018), who derived comparable
    judge/item-count curves in a different rating domain.
  \end{itemize}
\item
  \textbf{RQ2 (expertise).} Does rater expertise or exposure affect
  inter-rater agreement or rater severity? The
  Amabile-versus-Plucker/Kaufman tension developed in Sub-debate B above
  yields two competing, directional hypotheses rather than a single
  predicted winner; the RQ2 analyses test which position this dataset
  supports.

  \begin{itemize}
  \tightlist
  \item
    \textbf{H2a (Amabile-consistent).} Rater expertise level will not
    systematically affect agreement or severity.
  \item
    \textbf{H2b (Plucker/Kaufman-consistent).} Non-expert but
    high-exposure raters will show agreement or severity comparable to,
    or exceeding, that of formally credentialed experts.
  \end{itemize}
\item
  \textbf{RQ3 (aggregation/psychometrics).} Do Judge Response Theory
  and/or Many-Facet Rasch Model rater-severity models improve on simple
  mean aggregation, and what do they reveal about individual raters?
  This question is stated as open/exploratory, with no directional
  hypothesis: JRT, MFRM, and naive averaging have not previously been
  compared head-to-head at this scale in a creativity-rating context, so
  there is no established prior result from which to derive a
  directional prediction.
\item
  \textbf{RQ4 (construct validity).} What is the factor structure
  relating creativity, aesthetic appeal, and technical-correctness
  ratings? This question is likewise stated as open/exploratory: the
  factor structure among these three constructs is not yet established
  in this domain.
\item
  \textbf{RQ5 (reliability over time).} How long can raters perform the
  rating task reliably before agreement or rating quality declines? This
  question is stated as open/exploratory: no prior football-specific, or
  more broadly sport-specific, fatigue/session-length data exist to
  ground a directional prediction.
\end{itemize}

This paper addresses these five measurement questions only. Correlating
rated creativity with performance outcomes (e.g., the value of
possessions added metric, VAEP, or player market value) and developing
automated or machine-learning-based creativity scoring are both
explicitly deferred to two companion papers; the present paper is
intended to establish the measurement foundation on which those
applications will depend. Because RQ2 speaks directly to the practical
question of who should be recruited as a rater in future
football-creativity studies -- coaches and other domain experts, or
high-exposure non-expert viewers -- resolving it is intended as one of
this paper's principal applied contributions to sport science, a point
returned to in the Discussion and Conclusion.

\section{Methods}\label{methods}

This section describes a fully specified, reproducible proposed design.
Because data collection has not yet occurred, all procedures are
described in the future tense; nothing reported below has yet been
executed.

\subsection{Design overview}\label{design-overview}

This will be a cross-sectional, multi-rater creativity-rating study of
discrete football actions. The design combines a benchmark structure, in
which a fixed subset of actions is rated by every judge in the rater
panel, with a larger pairwise/rotating structure, in which a much larger
set of actions is rated by only two to three judges each. The benchmark
subset will enable full-overlap reliability analyses (RQ1-RQ3) at every
candidate judge count, while the pairwise subset will allow the dataset
to scale beyond what a fully crossed design could support within the
available rater capacity. The benchmark panel will comprise a fixed set
of \textbf{10 raters}, each of whom rates the full benchmark subset;
this yields ICC/agreement trajectories at every candidate judge count
from 1 through 10 (see ``Judges/raters'' and Analysis Plan (a) below). A
panel of this size is large enough that the study can directly test the
generic ``five to ten judges'' heuristic (Kaufman, Plucker, et al.,
2008) against its own benchmark data, rather than merely citing the
heuristic secondhand -- a genuine strength of this design, since the
resulting ICC-by-judge-count curve spans exactly the range the heuristic
claims is sufficient.

\subsection{Stimuli: Sample of football
actions}\label{stimuli-sample-of-football-actions}

Video and event data will be sourced preferentially from StatsBomb
(pending confirmation of video access) or, as a fallback, from WyScout
17/18. Matches will be sampled across multiple competitions, including
men's and women's FIFA tournaments, the Africa Cup of Nations, UEFA
tournaments, and the Copa America, in order to capture a broad range of
playing contexts and player populations.

The target sample is a purposive selection of approximately
\textbf{1,000 to 2,000 football actions}, pre-selected as plausible
creative-moment candidates rather than drawn through exhaustive or
random sampling of every action in every match. Candidate actions might
be identified via event-data heuristics (e.g., flags for unusual pass
trajectories, dribbles, or shot techniques) or through expert
pre-screening designed to surface actions likely to vary meaningfully in
creativity.

\subsubsection{Sample-size
justification}\label{sample-size-justification}

The \textasciitilde1,000-2,000-action target is, at present, a
plausibility-based design choice rather than the product of a formal a
priori sample-size or power calculation, and no such calculation has yet
been performed. A more rigorous alternative, should the author wish to
pursue it before data collection begins, would be precision-based
planning conducted within a generalizability-theory (G-theory)
framework: for each candidate number of judges under consideration for
the benchmark subset, the minimum benchmark size required to achieve a
target confidence-interval width for the resulting intraclass
correlation coefficient (ICC) could be estimated in advance, following
the logic used by Hass et al. (2018) to empirically derive the number of
raters and items needed for acceptable reliability of divergent-thinking
ratings. This is presented here as a decision point for the author:
either (a) proceed with the \textasciitilde1,000-2,000-action purposive
target as currently planned, treating it as a plausible and logistically
feasible starting point, or (b) conduct a formal precision- or
power-based calculation first and adjust the target accordingly. As of
this writing, option (a) is the operative plan, and this should be
understood as a provisional rather than a finalized decision.

\subsection{Judges/raters}\label{judgesraters}

Raters will be recruited to span an expertise gradient, including both
domain experts (e.g., coaches or other football professionals) and
non-expert but high-exposure viewers (e.g., committed spectators without
formal coaching or playing credentials). This recruitment plan directly
operationalizes the Amabile-versus-Plucker/Kaufman expertise debate
outlined in the Introduction (Amabile, 1982; Baer, 2015; Castel et al.,
2007; Ceh et al., 2021; Kaufman, Plucker, et al., 2008; Kaufman, Baer,
et al., 2008; Plucker et al., 2009): rather than assuming a single
``correct'' level of expertise, the design will allow expertise and
exposure effects on agreement and rater severity to be estimated
directly from the data.

The rater panel will be organized into two overlapping subsets. A fixed
benchmark panel of \textbf{10 raters} will rate every action in the
benchmark subset, enabling full pairwise ICC and agreement comparisons
at every candidate judge count from a single judge up to the full panel
of 10 (see Analysis Plan (a); the rationale for this panel size is given
in Design overview above). A larger pairwise/rotating subset of actions
will be rated by only two to three judges per action, drawn from the
benchmark panel and/or a separate pool of additional pairwise-only
raters; the exact size of this additional pool is not yet fixed and will
depend on recruitment logistics -- this is stated explicitly here as an
open item, not invented. The total rater pool for the study is therefore
the 10-rater benchmark panel plus this additional, as-yet-unsized
pairwise-only pool.

Raters assigned to the pairwise subset will be allocated to actions via
random assignment constrained to preserve a connected rater-by-action
network -- that is, assignment will not be fully unconstrained, but
constrained so that the resulting rater-by-action bipartite graph
remains connected, a standard requirement for JRT/MFRM parameter
identifiability. This constraint partially de-risks, but does not
eliminate, the rater-by-action sparsity concern discussed in Limitations
below: connectivity is guaranteed by design, but the realized degree of
overlap remains unknown until data collection occurs.

\subsection{Rating instrument}\label{rating-instrument}

Actions will be rated on a 7-point unipolar novelty x appropriateness
scale, with behaviorally anchored verbal descriptions at each scale
point, following psychometric guidance on selecting an appropriate
number of response categories. In addition to the primary creativity
rating, judges will also rate two secondary dimensions -- aesthetic
appeal and technical correctness -- for construct-validity purposes,
following Amabile's (1982) recommendation that secondary constructs be
rated alongside creativity whenever CAT is applied to a new domain, and
following the standard definitional consensus that creativity itself
reflects the joint presence of novelty and appropriateness (Runco \&
Jaeger, 2012). Judges will undergo a frame-of-reference (FOR) training
procedure prior to rating, intended to calibrate their use of the
scale's anchor points before live ratings begin: raters will review a
set of worked/anchor examples at each of the scale's seven points,
followed by a short calibration rating set with feedback, before
proceeding to live rating. This training session is expected to take
approximately 30-45 minutes, and raters will be required to reach an
agreement of at least .70 (e.g., intraclass correlation or an equivalent
percentage-agreement-within-one-scale-point criterion, with the exact
index to be finalized alongside the main analysis pipeline) with
criterion or consensus ratings on the calibration set before being
approved to proceed to live rating; raters who do not meet this
threshold will repeat calibration with an additional set of worked
examples. This calibration protocol is a plausible, generic design
choice consistent with standard FOR-training practice, and has not
itself been piloted. Raters will view full match footage and context,
including player and team identity, rather than blinded or anonymized
clips, prioritizing ecological validity -- creativity judgments grounded
in naturalistic match footage rather than standardized, blinded stimuli
-- over experimental blinding (see Limitations for the corresponding
halo-effect risk). An explicit conversion bridge will be specified
between this 7-point scale and the 0-10 scale used by Kempe and Memmert
(2018), to allow cross-study comparability with that benchmark study.

\subsection{Procedure}\label{procedure}

Ratings will be collected via a dedicated rating platform presenting
actions in randomized order and exporting structured rating data for
analysis; the specific platform has not yet been selected and is flagged
here, like the additional pairwise-rater-pool size above and the
pre-selection criteria in ``Stimuli,'' as an open item to be finalized
by the author before data collection. A pilot phase will precede full
data collection, in which per-session rating capacity will be
established by monitoring agreement and rating consistency as a function
of position within the rating session; this pilot directly informs RQ5
and follows the design logic of a prior study examining CAT reliability
over extended sessions (Barth \& Stadtmann, 2020).

\subsection{Analysis plan}\label{analysis-plan}

Six analyses are planned. Five map directly onto the five research
questions above; the sixth, (b), is an auxiliary/supporting analysis for
RQ2 rather than an independent research question of its own (see (b)
below). All six analyses (a)-(f) commit, as a matter of design, to
reporting confidence intervals and effect sizes alongside, or instead
of, significance tests.

\textbf{(a) Judge-count convergence (RQ1).} Intraclass correlation
coefficients (e.g., ICC(2,k) or ICC(3,k), as appropriate to the design)
will be computed on the benchmark subset and examined as a function of
the number of raters averaged into the aggregate rating (k = 1 through
the full 10-rater panel), with a possible Spearman-Brown-type
reliability projection to candidate judge counts beyond those directly
observed. A target reliability of ICC(2,k) \textgreater= .80 is
pre-specified as the criterion for declaring ratings
``stable/convergent'' for the purposes of RQ1, with ICC(2,k)
\textgreater= .70 noted as a minimum acceptable floor, following
standard convention in the reliability/psychometrics literature; the
number of judges ``needed'' will be read off the ICC-by-judge-count
curve against this pre-specified threshold, not chosen post hoc. This
analysis is motivated by Cseh and Jeffries' (2019) critique and follows
the empirical, G-theory-informed template set out by Hass et al. (2018),
providing a direct, domain-specific counterpart to the generic ``five to
ten judges'' heuristic (Kaufman, Plucker, et al., 2008).

\textbf{(b) Variance-source mixed model (auxiliary analysis for RQ2).} A
mixed-effects model of the form
\texttt{rating\ \textasciitilde{}\ covariates\ +\ (1\textbar{}rater)\ +\ (1\textbar{}player)\ +\ (1\textbar{}match)}
will be fit to partition variability in ratings across rater, player,
and match sources and to estimate the effect of any relevant covariates,
including rater expertise. This model is an auxiliary, supporting
analysis for RQ2 rather than an independent sixth research question: it
decomposes rating variance into rater/player/match sources and allows
expertise to be entered as a covariate or moderator, directly informing
the H2a/H2b comparison described in the Introduction.

\textbf{(c) Rater-severity diagnostics via JRT and MFRM (RQ3).} Rater
severity and leniency will be modeled using Judge Response Theory and/or
the Many-Facet Rasch Model, and the resulting rater-corrected estimates
will be compared against simple mean aggregation (Eckes, 2015; Hung et
al., 2012; Myford \& Wolfe, 2003, 2004; Myszkowski, 2021, 2024;
Myszkowski \& Storme, 2019, 2025; Primi et al., 2019).

\textbf{(d) Expertise-group comparison (RQ2).} Agreement statistics
(ICC) and/or rater-severity parameters from the JRT/MFRM models will be
compared across expertise groups; expertise will also be entered as a
covariate or moderator in the mixed model described above.

\textbf{(e) Factor analysis (RQ4).} An exploratory and/or confirmatory
factor analysis will be conducted across the creativity,
aesthetic-appeal, and technical-correctness ratings to evaluate their
dimensional structure (Amabile, 1982).

\textbf{(f) Fatigue/session-duration reliability (RQ5).} Agreement and
rating consistency will be examined as a function of session position or
time-on-task, using the pilot data described above, following the design
logic of Barth and Stadtmann (2020).

\subsubsection{Software and estimation}\label{software-and-estimation}

The full analysis pipeline will be implemented within a single
open-source, reproducible R environment. JRT models will be fit using
the \texttt{jrt} R package (Myszkowski, 2021). MFRM models will be fit
using an R-based implementation: the Test Analysis Modules package (TAM)
for a standard many-facet Rasch model, or, as an alternative, a Bayesian
many-facet Rasch model fit via \texttt{brms}, with rater,
player/examinee, and item specified as crossed random facets (Eckes,
2015; Myford \& Wolfe, 2003, 2004). Commercial Facets software will not
be used. This keeps JRT estimation, MFRM estimation, the mixed models,
and the factor analysis all reproducible from a single R codebase.

\subsection{Ethics statement}\label{ethics-statement}

Ethics approval for this project was obtained from the Ethics Committee
of the German Sport University Cologne (approval no. 248/2025). This
study was not preregistered.

\section{Results}\label{results}

No data have yet been collected, and no numerical results are reported
in this section. What follows is a structural description of the results
that will be reported once data collection and analysis are complete,
organized in the same order as the research questions above. Each
planned result is linked to the corresponding placeholder object.

\subsection{Sample description}\label{sample-description}

Once data collection is complete, Table~\ref{tbl-sample-description}
will report the achieved number of distinct football actions rated, the
number of matches and competitions sampled, the number of judges
recruited (against the 10-rater benchmark-panel target plus the
additional pairwise-only rater pool, whose size is still to be
finalized), and the number of ratings obtained per action in both the
benchmark and pairwise subsets, set against the
\textasciitilde1,000-2,000-action purposive target specified in Methods.

\begin{table}

\caption{\label{tbl-sample-description}Description of the collected
sample: number of distinct football actions rated, number of matches and
competitions sampled, number of judges/raters recruited (against the
10-rater benchmark-panel target plus the additional pairwise-only rater
pool, size still to be finalized), and ratings per action (benchmark
vs.~pairwise subset), reported against the planned purposive target of
\textasciitilde1,000-2,000 actions.}

\centering{

}

\end{table}%

\subsection{RQ1: Judge-count
convergence}\label{rq1-judge-count-convergence}

Figure~\ref{fig-icc-by-judge-count} will show inter-rater reliability
(ICC) as a function of the number of judges included in the aggregate
rating (k = 1 through 10), computed on the 10-rater benchmark subset
rated by all judges. Table~\ref{tbl-agreement-by-judge-count} will
tabulate the corresponding agreement and reliability statistics, along
with category-usage checks, across a set of candidate judge counts (2,
3, 5, and 10 judges, all achievable from the 10-rater benchmark panel).
Both will be read against the pre-specified ICC(2,k) \textgreater= .80
target and .70 minimum floor stated in Methods, Analysis Plan (a), to
identify the number of judges needed for stable, convergent ratings.
Together these will address RQ1 and H1, and provide the empirical
counterpart to the general ``five to ten judges'' heuristic discussed in
the Introduction -- a comparison the 10-rater benchmark-panel design was
built to support directly.

\begin{figure}

\centering{

}

\caption{\label{fig-icc-by-judge-count}Inter-rater reliability (ICC) as
a function of the number of judges included in the aggregate rating (k =
1-10), computed on the 10-rater benchmark subset, with reference lines
at the pre-specified ICC(2,k) \textgreater= .80 target and .70 minimum
floor.}

\end{figure}%

\begin{table}

\caption{\label{tbl-agreement-by-judge-count}Summary
agreement/reliability statistics (e.g., ICC with 95\% CIs and
category-usage checks) tabulated across candidate numbers of judges (2,
3, 5, and 10, all achievable from the 10-rater benchmark panel), with
the pre-specified ICC(2,k) \textgreater= .80 target / .70 floor noted
alongside the achieved values.}

\centering{

}

\end{table}%

\subsection{RQ2: Expertise effects}\label{rq2-expertise-effects}

Figure~\ref{fig-expertise-agreement} will compare rater agreement and
validity metrics across expertise groups (e.g., domain experts
vs.~non-expert, high-exposure raters), framed against the competing
H2a/H2b hypotheses stated in the Introduction, and
Table~\ref{tbl-expertise-effects} will report mixed-model estimates of
rater expertise as a predictor of rating level and/or variance, with
rater, player, and match modeled as random effects, per the
variance-source model specified in Methods (analysis (b) above, an
auxiliary analysis supporting RQ2). Whether the general
variance-decomposition mixed model (Methods, analysis (b)) is reported
as its own dedicated table in addition to
Table~\ref{tbl-expertise-effects}, or folded entirely into that table,
remains an open editorial decision for a future results-embedder pass
once the model has actually been fit.

\begin{figure}

\centering{

}

\caption{\label{fig-expertise-agreement}Rater agreement/validity metrics
compared across rater-expertise groups (e.g., expert
vs.~novice/high-viewership judges).}

\end{figure}%

\begin{table}

\caption{\label{tbl-expertise-effects}Mixed-model estimates of rater
expertise as a predictor of rating level/variance, with
rater/player/match random effects.}

\centering{

}

\end{table}%

\subsection{RQ3: Rater-severity and aggregation
(JRT/MFRM)}\label{rq3-rater-severity-and-aggregation-jrtmfrm}

Figure~\ref{fig-rater-severity-jrt} will present rater severity
estimates from the JRT and/or MFRM models (fit via the \texttt{jrt} R
package and TAM or \texttt{brms}, respectively), ordered by rater and
shown with uncertainty intervals. Table~\ref{tbl-rater-fit-diagnostics}
will report rater fit diagnostics (e.g., infit and outfit statistics, or
their JRT equivalents) flagging any raters with atypical response
patterns, and will present a comparison between the rater-corrected
estimates and simple mean aggregation, addressing RQ3.

\begin{figure}

\centering{

}

\caption{\label{fig-rater-severity-jrt}Judge Response Theory (via the R
\texttt{jrt} package) / Many-Facet Rasch Model (via TAM or brms) rater
severity estimates, ordered by rater, with uncertainty intervals.}

\end{figure}%

\begin{table}

\caption{\label{tbl-rater-fit-diagnostics}Rater fit diagnostics (e.g.,
infit/outfit statistics or JRT-equivalent) flagging raters with atypical
response patterns.}

\centering{

}

\end{table}%

\subsection{RQ4: Construct and factor
validity}\label{rq4-construct-and-factor-validity}

Figure~\ref{fig-factor-structure} will present the factor or correlation
structure relating creativity (novelty x appropriateness),
aesthetic-appeal, and technical-correctness ratings, and
Table~\ref{tbl-factor-loadings} will report the corresponding factor
loadings and model fit indices, addressing RQ4.

\begin{figure}

\centering{

}

\caption{\label{fig-factor-structure}Factor structure / correlation
structure among creativity (novelty x appropriateness), aesthetic, and
technical-correctness ratings.}

\end{figure}%

\begin{table}

\caption{\label{tbl-factor-loadings}Factor loadings and fit indices for
the proposed factor model relating creativity, aesthetic appeal, and
technical correctness dimensions.}

\centering{

}

\end{table}%

\subsection{RQ5: Reliability over session
duration}\label{rq5-reliability-over-session-duration}

Findings addressing RQ5 -- whether agreement or rating consistency
declines as a function of session position or time-on-task, based on the
pilot fatigue-monitoring procedure described in Methods -- will be
reported narratively in this subsection once the pilot data are
available. No dedicated figure or table placeholder currently exists for
this analysis (see Limitations); the finding will therefore be described
in prose only, and/or folded into the sample-description reporting
above, pending a future results-embedder pass once real pilot data
exist.

\section{Discussion}\label{discussion}

\subsection{Contribution and place in the
literature}\label{contribution-and-place-in-the-literature}

This paper is intended as a direct empirical response to Cseh and
Jeffries' (2019) call for domain-specific evidence on CAT's design
parameters, and as an applied test, at scale, of the rater-effects
modeling program initiated by Myszkowski and Storme (2019). Whatever the
eventual results, generating a large-scale, systematically varied
multi-rater creativity-rating dataset for football actions is itself a
contribution: prior sport-creativity research has never varied judge
number or expertise systematically, nor compared rater-effects models
against naive averaging in this domain.

\textbf{{[}TBC -- author to confirm once data are available{]}} The
judge-count finding (RQ1) will bear on what number of judges should be
recommended for football, and by extension similarly naturalistic,
product-based creativity studies, read against the pre-specified
ICC(2,k) \textgreater= .80 target (.70 minimum floor; Methods, Analysis
Plan (a)). This eventual recommendation should be explicitly compared
against the generic ``five to ten judges'' heuristic (Kaufman, Plucker,
et al., 2008) (noting that the precise page-level source of this figure
remains unverified) and against the generalizability-theory-derived
numbers found in other rating domains (Hass et al., 2018) -- a
comparison the 10-rater benchmark-panel design was built to support
directly. Whether the football-specific number converges with, falls
below, or exceeds these comparators is not yet known.

\textbf{{[}TBC -- author to confirm once data are available{]}} The
expertise finding (RQ2) will indicate whether football creativity rating
resolves toward hypothesis H2a, Amabile's (1982) position that expertise
need not be uniform across judges, or toward hypothesis H2b, the
position that domain expertise does not straightforwardly improve
agreement (Kaufman, Baer, et al., 2008; Plucker et al., 2009). The
practical question of who should be recruited as a rater in future
football-creativity studies -- coaches and other domain experts, or
high-exposure non-expert viewers -- depends directly on this result.

\textbf{{[}TBC -- author to confirm once data are available{]}} The
rater-severity finding (RQ3) will speak to whether JRT and/or MFRM
modeling is worth the added analytic complexity relative to naive
averaging for this kind of naturalistic, non-parallel rating data. This
positions the present study as an applied test of the JRT research
program (Myszkowski, 2021, 2024; Myszkowski \& Storme, 2025) and of MFRM
as an established alternative (Eckes, 2015; Hung et al., 2012; Primi et
al., 2019), with both estimated throughout in R (\texttt{jrt}; TAM or
\texttt{brms}) for full reproducibility.

\textbf{{[}TBC -- author to confirm once data are available{]}} The
factor-structure finding (RQ4) will indicate whether creativity is
empirically discriminable from aesthetic appeal and technical
correctness in this domain, connecting back to the
taxonomy/facet-clarity argument that different facets of creativity
require distinct measurement treatment (Barbot et al., 2019) and to
Amabile's (1982) construct-validity rationale for rating secondary
dimensions alongside creativity.

\textbf{{[}TBC -- author to confirm once data are available{]}} The
fatigue/ session-length finding (RQ5) will provide practical guidance
for designing CAT rating protocols generally, not only in sport,
extending the precedent set by Barth and Stadtmann (2020).

\subsection{Generalizability}\label{generalizability}

Taken together, the eventual findings are intended to inform CAT
practice broadly, directly answering Cseh and Jeffries' (2019) call for
domain-specific empirical evidence on judge number, expertise, and
aggregation method. At the same time, a single applied domain --
football -- cannot by itself resolve a domain-general question; the
extent to which these findings transfer to other creative-product
domains remains an open question for future replication.

\subsection{Limitations}\label{limitations}

\begin{itemize}
\tightlist
\item
  This study addresses a single applied domain (football);
  generalization of any findings to other creative-product domains is an
  open question.
\item
  Rater recruitment is convenience- and plan-dependent, and the exact
  operationalization of ``expertise'' (e.g., which credentials or
  exposure thresholds distinguish expert from non-expert raters) still
  needs to be finalized before data collection.
\item
  The ``five to ten judges'' heuristic's exact source page has not been
  independently verified; it should not be cited as a precisely sourced
  figure until confirmed.
\item
  The identifiability of the MFRM/JRT models depends on the realized
  sparsity of the rater-by-action rating matrix, which is not yet known
  and cannot be guaranteed in advance of data collection; the
  constrained random-assignment rule described in Methods
  (``Judges/raters'') guarantees a connected rater-by-action network by
  design, which partially de-risks but does not eliminate this concern
  -- the realized degree of overlap remains unknown until data
  collection occurs.
\item
  Raters will view full match footage and context, including player and
  team identity, rather than blinded or anonymized clips (Methods,
  ``Rating instrument''), prioritizing ecological validity over
  blinding; this introduces a plausible halo-effect risk, in which
  recognizable players or teams could influence creativity judgments
  independently of the action itself. The Many-Facet Rasch Model
  (Sub-debate C; Analysis Plan (c)) is explicitly intended to help
  diagnose such rater idiosyncrasies, including halo effects, but cannot
  fully substitute for blinding.
\item
  Rater fatigue and motivation effects in a large-scale, potentially
  online/remote rating task are a genuine risk; no dedicated Results
  artifact is currently planned for the fatigue/session-length finding
  (RQ5), and a dedicated fatigue-analysis output still needs to be
  finalized in a future results-embedder pass once pilot data exist.
\item
  The purposive (non-random) sampling of ``plausible creative-moment''
  actions means that prevalence or base-rate claims about creativity in
  football cannot be supported by this sample; only relative,
  psychometric comparisons (judge-count, expertise, aggregation method)
  are licensed by this design.
\item
  No formal a priori sample-size or power calculation was performed for
  the \textasciitilde1,000-2,000-action target; it reflects a
  plausibility-based design decision rather than a precision-targeted
  calculation (see Methods).
\item
  This study was not preregistered.
\end{itemize}

\section{Conclusion}\label{conclusion}

This paper proposes an empirically grounded, reproducible methodological
template for consensual-assessment-technique-based creative-product
rating in sport, directly addressing a domain-specificity gap that the
broader creativity-assessment literature has repeatedly flagged but
rarely operationalized (Cseh \& Jeffries, 2019).

\textbf{{[}TBC -- author to confirm once data are available{]}} Once the
planned analyses are complete, this paper will offer practical
recommendations for future sport-creativity CAT studies concerning the
number of judges to recruit, the expertise mix of the rater panel, and
whether rater-effects modeling (JRT/MFRM) is worth adopting over simple
averaging. Resolving the expertise question (RQ2) in particular is
intended as one of the paper's principal applied contributions to sport
science, since it bears directly on who should be recruited as a rater
in future studies. These recommendations are contingent on the Results
and Discussion above and are not stated as specific numbers here.

This paper is the first of a planned three-paper sequence. A companion
application paper will examine the relationship between rated creativity
and performance outcomes such as the value of possessions added (VAEP)
and player market value; a further companion paper will explore
automated, machine-learning-based creativity scoring, using the human
ratings generated here as ground truth -- following a precedent, in the
same target journal, for validating automated scoring against
consensual-assessment-technique ratings of creative products (DiStefano
et al., 2024). Both are named here only as future work and are not
elaborated further. Getting the human measurement foundation right is
treated here as a precondition for both of these follow-on efforts.

\section{References}\label{references}

\phantomsection\label{refs}
\begin{CSLReferences}{1}{0}
\bibitem[\citeproctext]{ref-amabile1982}
Amabile, T. M. (1982). Social psychology of creativity: A consensual
assessment technique. \emph{Journal of Personality and Social
Psychology}, \emph{43}(5), 997--1013.
\url{https://doi.org/10.1037/0022-3514.43.5.997}

\bibitem[\citeproctext]{ref-baer2015}
Baer, J. (2015). The importance of domain-specific expertise in
creativity. \emph{Roeper Review}, \emph{37}(3), 165--178.
\url{https://doi.org/10.1080/02783193.2015.1047480}

\bibitem[\citeproctext]{ref-baer2004}
Baer, J., Kaufman, J. C., \& Gentile, C. A. (2004). Extension of the
consensual assessment technique to nonparallel creative products.
\emph{Creativity Research Journal}, \emph{16}(1), 113--117.
\url{https://doi.org/10.1207/s15326934crj1601_11}

\bibitem[\citeproctext]{ref-barbot2019b}
Barbot, B., Hass, R. W., \& Reiter-Palmon, R. (2019). Creativity
assessment in psychological research: (Re)setting the standards.
\emph{Psychology of Aesthetics, Creativity, and the Arts}, \emph{13}(2),
233--240. \url{https://doi.org/10.1037/aca0000233}

\bibitem[\citeproctext]{ref-barbot2019}
Barbot, B., \& Reiter-Palmon, R. (2019). Creativity assessment:
Pitfalls, solutions, and standards. \emph{Psychology of Aesthetics,
Creativity, and the Arts}, \emph{13}(2), 131--132.
\url{https://doi.org/10.1037/aca0000251}

\bibitem[\citeproctext]{ref-barth2020}
Barth, P., \& Stadtmann, G. (2020). Creativity assessment over time:
Examining the reliability of {CAT} ratings. \emph{The Journal of
Creative Behavior}, \emph{55}(2), 396--409.
\url{https://doi.org/10.1002/jocb.462}

\bibitem[\citeproctext]{ref-batey2012}
Batey, M. (2012). The measurement of creativity: From definitional
consensus to the introduction of a new heuristic framework.
\emph{Creativity Research Journal}, \emph{24}(1), 55--65.
\url{https://doi.org/10.1080/10400419.2012.649181}

\bibitem[\citeproctext]{ref-castel2007}
Castel, A. D., McCabe, D. P., Roediger, H. L., \& Heitman, J. L. (2007).
The dark side of expertise. \emph{Psychological Science}, \emph{18}(1),
3--5. \url{https://doi.org/10.1111/j.1467-9280.2007.01838.x}

\bibitem[\citeproctext]{ref-ceh2021}
Ceh, S. M., Edelmann, C., Hofer, G., \& Benedek, M. (2021). Assessing
raters: What factors predict discernment in novice creativity raters?
\emph{The Journal of Creative Behavior}, \emph{56}(1), 41--54.
\url{https://doi.org/10.1002/jocb.515}

\bibitem[\citeproctext]{ref-christiaans2005}
Christiaans, H., \& Venselaar, K. (2005). Creativity in design
engineering and the role of knowledge: Modelling the expert.
\emph{International Journal of Technology and Design Education},
\emph{15}(3), 217--236. \url{https://doi.org/10.1007/s10798-004-1904-4}

\bibitem[\citeproctext]{ref-cseh2019}
Cseh, G. M., \& Jeffries, K. K. (2019). A scattered {CAT}: A critical
evaluation of the consensual assessment technique for creativity
research. \emph{Psychology of Aesthetics, Creativity, and the Arts},
\emph{13}(2), 159--166. \url{https://doi.org/10.1037/aca0000220}

\bibitem[\citeproctext]{ref-distefano2024}
DiStefano, P. V., Patterson, J. D., \& Beaty, R. E. (2024). Automatic
scoring of metaphor creativity with large language models.
\emph{Creativity Research Journal}, 1--15.
\url{https://doi.org/10.1080/10400419.2024.2326343}

\bibitem[\citeproctext]{ref-eckes2015}
Eckes, T. (2015). \emph{Introduction to many-facet rasch measurement:
Analyzing and evaluating rater-mediated assessments} (2nd revised and
updated, Vol. 22). Peter Lang.
\url{https://doi.org/10.3726/978-3-653-04844-5}

\bibitem[\citeproctext]{ref-hass2018}
Hass, R. W., Rivera, M., \& Silvia, P. J. (2018). On the dependability
and feasibility of layperson ratings of divergent thinking.
\emph{Frontiers in Psychology}, \emph{9}, 1343.
\url{https://doi.org/10.3389/fpsyg.2018.01343}

\bibitem[\citeproctext]{ref-hung2012}
Hung, S.-P., Chen, P.-H., \& Chen, H.-C. (2012). Improving creativity
performance assessment: A rater effect examination with many facet
{Rasch} model. \emph{Creativity Research Journal}, \emph{24}(4),
345--357. \url{https://doi.org/10.1080/10400419.2012.730331}

\bibitem[\citeproctext]{ref-kaufman2008b}
Kaufman, J. C., Baer, J., Cole, J. C., \& Sexton, J. D. (2008). A
comparison of expert and nonexpert raters using the consensual
assessment technique. \emph{Creativity Research Journal}, \emph{20}(2),
171--178. \url{https://doi.org/10.1080/10400410802059929}

\bibitem[\citeproctext]{ref-kaufman2008}
Kaufman, J. C., Plucker, J. A., \& Baer, J. (2008). \emph{Essentials of
creativity assessment}. Wiley.

\bibitem[\citeproctext]{ref-kempe2018}
Kempe, M., \& Memmert, D. (2018). {``}Good, better, creative{''}: The
influence of creativity on goal scoring in elite soccer. \emph{Journal
of Sports Sciences}, \emph{36}(21), 2419--2423.
\url{https://doi.org/10.1080/02640414.2018.1459153}

\bibitem[\citeproctext]{ref-long2022}
Long, H., Kerr, B. A., Emler, T. E., \& Birdnow, M. (2022). A critical
review of assessments of creativity in education. \emph{Review of
Research in Education}, \emph{46}(1), 288--323.
\url{https://doi.org/10.3102/0091732X221084326}

\bibitem[\citeproctext]{ref-long2015}
Long, H., \& Pang, W. (2015). Rater effects in creativity assessment: A
mixed methods investigation. \emph{Thinking Skills and Creativity},
\emph{15}, 13--25. \url{https://doi.org/10.1016/j.tsc.2014.10.004}

\bibitem[\citeproctext]{ref-memmert2011}
Memmert, D. (2010). Creativity, expertise, and attention: Exploring
their development and their relationships. \emph{Journal of Sports
Sciences}, \emph{29}(1), 93--102.
\url{https://doi.org/10.1080/02640414.2010.528014}

\bibitem[\citeproctext]{ref-myford2003}
Myford, C. M., \& Wolfe, E. W. (2003).
\href{https://www.ncbi.nlm.nih.gov/pubmed/14523257}{Detecting and
measuring rater effects using many-facet {Rasch} measurement: {P}art
{I}}. \emph{Journal of Applied Measurement}, \emph{4}(4), 386--422.

\bibitem[\citeproctext]{ref-myford2004}
Myford, C. M., \& Wolfe, E. W. (2004).
\href{https://www.ncbi.nlm.nih.gov/pubmed/15064538}{Detecting and
measuring rater effects using many-facet {Rasch} measurement: {P}art
{II}}. \emph{Journal of Applied Measurement}, \emph{5}(2), 189--227.

\bibitem[\citeproctext]{ref-myszkowski2021}
Myszkowski, N. (2021). Development of the {R} library "jrt": Automated
item response theory procedures for judgment data and their application
with the consensual assessment technique. \emph{Psychology of
Aesthetics, Creativity, and the Arts}, \emph{15}(3), 426--438.
\url{https://doi.org/10.1037/aca0000287}

\bibitem[\citeproctext]{ref-myszkowski2024}
Myszkowski, N. (2024). \emph{Item response theory for creativity
measurement}. Cambridge University Press.
\url{https://doi.org/10.1017/9781009239035}

\bibitem[\citeproctext]{ref-myszkowski2019}
Myszkowski, N., \& Storme, M. (2019). Judge response theory? A call to
upgrade our psychometrical account of creativity judgments.
\emph{Psychology of Aesthetics, Creativity, and the Arts}, \emph{13}(2),
167--175. \url{https://doi.org/10.1037/aca0000225}

\bibitem[\citeproctext]{ref-myszkowski2025}
Myszkowski, N., \& Storme, M. (2025). One score, two components:
Disentangling appropriateness and originality in {PISA} creative
thinking judgments using generalized item response tree models.
\emph{The Journal of Creative Behavior}, \emph{59}(2), e70033.

\bibitem[\citeproctext]{ref-oboeuf2020}
Oboeuf, A., Hanneton, S., Buffet, J., Fantoni, C., \& Labiadh, L.
(2020). Influence of traditional sporting games on the development of
creative skills in team sports. The case of football. \emph{Frontiers in
Psychology}, \emph{11}, 611803.
\url{https://doi.org/10.3389/fpsyg.2020.611803}

\bibitem[\citeproctext]{ref-plucker2009}
Plucker, J. A., Kaufman, J. C., Temple, J. S., \& Qian, M. (2009). Do
experts and novices evaluate movies the same way? \emph{Psychology \&
Marketing}, \emph{26}(5), 470--478.
\url{https://doi.org/10.1002/mar.20283}

\bibitem[\citeproctext]{ref-primi2019}
Primi, R., Silvia, P. J., Jauk, E., \& Benedek, M. (2019). Applying
many-facet rasch modeling in the assessment of creativity.
\emph{Psychology of Aesthetics, Creativity, and the Arts}, \emph{13}(2),
176--186. \url{https://doi.org/10.1037/aca0000230}

\bibitem[\citeproctext]{ref-runco2012}
Runco, M. A., \& Jaeger, G. J. (2012). The standard definition of
creativity. \emph{Creativity Research Journal}, \emph{24}(1), 92--96.
\url{https://doi.org/10.1080/10400419.2012.650092}

\bibitem[\citeproctext]{ref-safardilha2019}
Sa Fardilha, F. de, \& Allen, J. B. (2019). Defining, assessing, and
developing creativity in sport: A systematic narrative review.
\emph{International Review of Sport and Exercise Psychology},
\emph{13}(1), 104--127.
\url{https://doi.org/10.1080/1750984x.2019.1616315}

\bibitem[\citeproctext]{ref-zahno2022}
Zahno, S., \& Hossner, E.-J. (2022). Creative actions in team sports are
rooted in motor skills rather than in a divergent thinking ability.
\emph{German Journal of Exercise and Sport Research}, \emph{53}(2),
206--216. \url{https://doi.org/10.1007/s12662-022-00847-6}

\end{CSLReferences}




\end{document}
